% chapters/18_simulation_results.tex
% بخش پنجم: نتایج شبیه‌سازی‌ها
% این فایل شامل تمام نمودارهای خروجی ۱۰ کمپین شبیه‌سازی است.
% پیش‌نیاز: بسته‌های زیر باید در preamble فایل اصلی اضافه شوند:
%   \usepackage{subcaption}   % برای subfigure ها
%   \usepackage{float}         % برای گزینه [H] در صورت نیاز
% مسیر تصاویر: فرض شده PNG ها در پوشه figures/ کنار فایل اصلی قرار دارند.

% --- helper: include image with graceful fallback if file is missing ---
% \detokenize renders the filename as plain catcode-12 characters so that
% underscores (and other special tokens) inside the path do NOT trigger
% "Missing $ inserted" / "Extra }" / "There's no line here to end" errors
% when the file is absent and the fallback box is typeset.
\providecommand{\safeimg}[2][]{%
\IfFileExists{#2}{\includegraphics[#1]{#2}}{%
	\fbox{\parbox[c][5cm][c]{0.8\textwidth}{\centering
			\ttfamily\detokenize{#2}\par\vspace{4pt}%
			\normalfont\small(فایل تصویری یافت نشد --- تصویر را در مسیر فوق قرار دهید)%
	}}%
}%
}

% =====================================================================
% PART V: نتایج شبیه‌سازی
% =====================================================================
\part{نتایج شبیه‌سازی}
\label{part:results}

% ---------------------------------------------------------------------
\chapter{تنظیمات آزمایش و کمپین‌های شبیه‌سازی}
\label{ch:results_overview}

\section{نمای کلی کمپین‌ها}
این بخش نتایج حاصل از اجرای ده کمپین شبیه‌سازی مستقل با اسکریپت \lr{\texttt{main\_optimized.py}} را گزارش می‌کند. هر کمپین با ترکیب متفاوتی از پروتکل، اثبات امنیت، نوع آشکارساز و پارامترهای فیزیکی اجرا شده تا رفتار چارچوب در شرایط عملیاتی گوناگون ارزیابی شود. خروجی هر اجرا در یک فایل \lr{CSV} به نام \lr{\texttt{qkd\_results\_optimized\_sweep.csv}} نوشته می‌شود و سپس برای تحلیل به نمودار \lr{PNG} تبدیل شده است. این نمودارها در ادامهٔ فصل به‌ترتیب موضوعی ارائه می‌شوند.

تنها خروجی ملموس هر شبیه‌سازی همین فایل \lr{CSV} است؛ اسکریپت هیچ نمودار، فایل \lr{JSON} خلاصه یا گزارش تجمیعی تولید نمی‌کند. بنابراین تمام تحلیل‌های آماری بعدی بر مبنای پردازش همین فایل با کتابخانهٔ \lr{\texttt{pandas}} انجام شده است. تعداد کل شبیه‌سازی‌های اجرا شده در این ده کمپین تقریباً برابر ۱۸۰۰ نقطه است که هر کدام شامل تولید پالس، انتقال در کانال فیبر، آشکارسازی، غربال‌سازی پایه‌ها (\lr{sifting})، تخمین پارامتر و محاسبهٔ کلید امن هستند.

\section{خلاصهٔ کمپین‌ها}
جدول~\ref{tab:campaigns_summary} ده کمپین اجرا شده را به‌همراه پارامترهای متمایزکننده و تعداد نقاط شبیه‌سازی خلاصه می‌کند. این جدول به‌عنوان مرجع سریع برای تفسیر نمودارهای فصل‌های بعدی استفاده می‌شود.

\begin{table}[htbp]
\centering
\caption{خلاصهٔ کمپین‌های شبیه‌سازی اجرا شده با \lr{\texttt{main\_optimized.py}}}
\label{tab:campaigns_summary}
\small
\begin{tabular}{@{}cllc@{}}
	\toprule
	\textbf{شماره} & \textbf{نام فایل} & \textbf{محور تغییر اصلی} & \textbf{تعداد نقاط} \\
	\midrule
	۱ & \lr{\texttt{01\_baseline\_BB84\_Lim2014}}        & خط پایه \lr{BB84} با \lr{Lim2014}            & ۹۶  \\
	۲ & \lr{\texttt{02\_finite\_key\_scaling}}            & مقیاس پالس $10^4$ تا $10^8$       & ۲۸۸ \\
	۳ & \lr{\texttt{03\_dwdm\_*}}                          & \lr{DWDM} با توان‌های $-20$، $-34$ و $-40$ \lr{dBm} & ۲۸۸ \\
	۴ & \lr{\texttt{04\_auto\_optimized\_pulses}}        & بهینه‌سازی \lr{SLSQP} پارامتر پالس       & ۹۶  \\
	۵ & \lr{\texttt{05\_protocol\_*}}                      & چهار پروتکل مختلف                   & ۳۸۴ \\
	۶ & \lr{\texttt{06\_proof\_*}}                         & چهار اثبات امنیت متفاوت             & ۳۸۴ \\
	۷ & \lr{\texttt{07\_detector\_*}}                      & سه نوع آشکارساز \lr{SPD}/\lr{SNSPD}/\lr{PNRD}     & ۲۸۸ \\
	۸ & \lr{\texttt{08\_confidence\_*}}                    & سه روش بازهٔ اطمینان                 & ۲۸۸ \\
	۹ & \lr{\texttt{09\_rogers\_validation\_dense}}      & اعتبارسنجی \lr{Rogers} با ۵۱ نقطه        & ۱۰۲ \\
	۱۰ & \lr{\texttt{10\_qber\_misalignment\_sweep}}      & حساسیت به \lr{QBER} و \lr{misalignment}       & ۱۹۸ \\
	\bottomrule
\end{tabular}
\end{table}

\section{پیکربندی مشترک}
تمام کمپین‌ها از مقادیر پیش‌فرض زیر برای اطمینان از قابلیت مقایسه استفاده می‌کنند: افت فیبر \lr{0.2 dB/km}، بازهٔ فاصلهٔ \lr{0} تا \lr{150 km} با ۱۶ نقطه، دانهٔ تصادفی (\lr{seed}) برابر \lr{12345}، و هشت پردازشگر موازی. تعداد پالس‌ها مگر در کمپین ۲، در بازهٔ $10^4$ تا $10^6$ نگه داشته شده است. هر نقطه در دو حالت \lr{\texttt{noise\_applied=True}} و \lr{\texttt{noise\_applied=False}} اجرا می‌شود تا تأثیر نویز واقعی آشکارساز (شامل شمارش تاریک، پس‌پالس، جیتر، زمان مرده و عدم انطباق پایه‌ها) به‌صراحت قابل ارزیابی باشد.

% ---------------------------------------------------------------------
\chapter{نتایج خط پایه و مقیاس کلید محدود}
\label{ch:results_baseline}

\section{خط پایه BB84 با اثبات Lim2014}
شکل~\ref{fig:01_baseline} منحنی نرخ کلید امن به‌عنوان تابعی از فاصلهٔ فیبر را برای پروتکل \lr{BB84-Decoy} با اثبات \lr{Lim2014} نشان می‌دهد. این منحنی به‌عنوان مرجع اصلی گزارش عمل می‌کند و همهٔ مقایسه‌های بعدی نسبت به آن تفسیر می‌شوند. در حالت بدون نویز، نرخ کلید در فواصل کوتاه به مقدار نظری نزدیک می‌شود و با افزایش فاصله به‌دلیل افت نمایی فیبر ($\eta \propto 10^{-\alpha L/10}$) کاهش می‌یابد. در حالت با نویز واقعی، افت اضافی ناشی از شمارش تاریک (\lr{dark count}) و عدم انطباق پایه‌ها (\lr{misalignment}) باعث می‌شود نرخ کلید در فواصل طولانی‌تر سریع‌تر به صفر برسد.

\begin{figure}[htbp]
\centering
\safeimg[width=0.85\textwidth]{figures/01_baseline_BB84_Lim2014.csv.png}
\caption{نرخ کلید امن به‌ازای پالس در برابر فاصلهٔ فیبر برای پروتکل \lr{BB84-Decoy} با اثبات \lr{Lim2014}. دو منحنی حالت بدون نویز و با نویز واقعی آشکارساز مقایسه شده‌اند.}
\label{fig:01_baseline}
\end{figure}

نکتهٔ کلیدی در این شکل، شکستگی منحنی در فاصله‌ای است که در آن نرخ کلید به صفر می‌رسد؛ این نقطه نشان‌دهندهٔ حداکثر برد امن پروتکل در شرایط داده‌شده است. مقادیر دقیق این برد برای ترکیب‌های مختلف پارامتر در فصل‌های بعدی مقایسه می‌شود. علاوه بر این، ستون \lr{\texttt{status}} در \lr{CSV} خروجی، نقاطی که به وضعیت \lr{\texttt{ZERO\_KEY}} ختم شده‌اند را علامت‌گذاری می‌کند که برای تعیین برد امن معتبر هستند.

\section{مقیاس کلید محدود (\lr{Finite-Key Scaling})}
شکل~\ref{fig:02_finite_key} اثر اندازهٔ کلید بر نرخ کلید امن را در پنج اندازهٔ پالس از $10^4$ تا $10^7$ نشان می‌دهد. در اندازه‌های کوچک، جریمهٔ کلید محدود (\lr{finite-size penalty}) نرخ امن را به‌طور قابل توجهی کاهش می‌دهد؛ زیرا تخمین پارامتر با تعداد محدودی نمونه انجام می‌شود و بازهٔ اطمینان \lr{QBER} گسترده‌تر است. با افزایش تعداد پالس به $10^7$، نرخ کلید به مقدار مجانبی نزدیک می‌شود.

\begin{figure}[htbp]
\centering
\safeimg[width=0.85\textwidth]{figures/02_finite_key_scaling.csv.png}
\caption{اثر تعداد پالس بر نرخ کلید امن. با افزایش $N$ از $10^4$ به $10^7$، جریمهٔ کلید محدود کاهش می‌یابد و نرخ به مقدار مجانبی همگرا می‌شود.}
\label{fig:02_finite_key}
\end{figure}

این نتیجه اهمیت انتخاب اندازهٔ بلوک مناسب در پیاده‌سازی عملی \lr{QKD} را تأیید می‌کند؛ برای کاربردهای بلندبرد، اندازهٔ بلوک حداقل $10^6$ پالس ضروری است تا جریمهٔ \lr{finite-size} چشم‌پوشی‌پذیر شود. در مقابل، برای کاربردهای کوتاه‌برد با نرخ بالای رویداد، حتی $10^5$ پالس ممکن است کافی باشد.

\section{اعتبارسنجی Rogers 2007}
شکل~\ref{fig:09_rogers} مقایسهٔ نرخ کلید مونت‌کارلو (\lr{Monte Carlo}) به‌دست آمده از شبیه‌سازی کامل را با فرمول تحلیلی \lr{Rogers et al. 2007} (رابطهٔ ۱۵) برای نرخ بیت غربال‌شده (\lr{sifted bit}) نشان می‌دهد. این مقایسه در ۵۱ نقطهٔ متراکم از ۰ تا ۱۰۰ کیلومتر انجام شده تا هم‌خوانی شبیه‌سازی با تئوری به‌خوبی بررسی شود. ستون \lr{\texttt{analytic\_rogers\_sbr}} در فایل \lr{CSV} به‌طور خودکار برای هر نقطه محاسبه می‌شود و این امکان را فراهم می‌کند که نمودار دقیقاً بازتولید اشکال ۳ تا ۵ مقالهٔ \lr{Rogers} باشد.

\begin{figure}[htbp]
\centering
\safeimg[width=0.85\textwidth]{figures/09_rogers_validation_dense.csv.png}
\caption{مقایسهٔ نرخ \lr{sifted bit} تحلیلی \lr{Rogers 2007} با نرخ کلید امن مونت‌کارلو. هم‌خوانی دو منحنی، اعتبارسنجی صحت پیاده‌سازی لایهٔ فیزیکی و آشکارسازی را تأیید می‌کند.}
\label{fig:09_rogers}
\end{figure}


% ---------------------------------------------------------------------
\chapter{سناریوهای DWDM و بهینه‌سازی پالس}
\label{ch:results_dwdm_opt}

\section{تأثیر هم‌زیستی با کانال کلاسیک (\lr{DWDM})}
شکل~\ref{fig:03_dwdm_power} سه نمودار کنار هم را برای توان‌های گیرندهٔ مختلف \lr{$-20$}، \lr{$-34$} و \lr{$-40$ dBm} نشان می‌دهد و اثر نویز رامان بر نرخ کلید امن را روشن می‌سازد. در توان بالاتر (\lr{$-20$ dBm})، نویز رامان نرخ تیرگی مؤثر آشکارساز را به‌طور چشمگیری افزایش می‌دهد و حداکثر برد امن را به‌میزان قابل توجهی کاهش می‌دهد. در توان پایین‌تر (\lr{$-40$ dBm})، رفتار به حالت فیبر اختصاصی نزدیک می‌شود و افت فیلتر \lr{AWG} (۳ دسی‌بل) تنها منبع انحراف باقی می‌ماند.


\begin{figure}[htbp]
	\centering
	\safeimg[width=0.85\textwidth]{figures/03_dwdm_4channel.csv.png}
	\caption{توان $-34$ dBm (پیش‌فرض Lim2014)}
	\label{fig:03a_dwdm_default}
\end{figure}
\begin{figure}[htbp]
	\centering
	\safeimg[width=0.85\textwidth]{figures/03b_dwdm_power_minus20.csv.png}
	\caption{توان $-20$ dBm (نویز شدید)}
	\label{fig:03b_dwdm_high}
\end{figure}
\begin{figure}[htbp]
	\centering
	\safeimg[width=0.85\textwidth]{figures/03c_dwdm_power_minus40.csv.png}
	\caption{توان $-40$ dBm (نویز کم)}
	\label{fig:03c_dwdm_low}
\end{figure}


این نتایج با بخش IV مقالهٔ \lr{Lim 2014} همخوانی دارد و نشان می‌دهد که طراحی سیستم \lr{DWDM} برای \lr{QKD} نیازمند موازنهٔ دقیق بین توان کانال کلاسیک (برای ارتباط با نرخ بالا) و نویز رامان تحمیل‌شده بر کانال کوانتومی است. مقدار پیش‌فرض \lr{$-34$ dBm} نقطهٔ تعادل پیشنهادی مقاله است که در این شبیه‌سازی بازتولید شده است.

\section{بهینه‌سازی خودکار پارامتر پالس}
شکل~\ref{fig:04_auto_opt} نتیجهٔ بهینه‌سازی \lr{SLSQP} برای پنج پارامتر پالس $(\mu, \nu, p_s, p_d, p_v)$ به‌ازای هر فاصله را نشان می‌دهد. بهینه‌ساز از نرخ \lr{GLLP} مجانبی با جریمهٔ کلید محدود به‌عنوان تابع هدف استفاده می‌کند و شدت سیگنال $\mu$ را در فواصل کوتاه بالا و در فواصل طولانی پایین نگه می‌دارد تا تعادل بین بازدهی آشکارسازی و افشای اطلاعات بهینه شود.

\begin{figure}[htbp]
\centering
\safeimg[width=0.85\textwidth]{figures/04_auto_optimized_pulses.csv.png}
\caption{نرخ کلید امن با بهینه‌سازی خودکار پارامتر پالس در برابر خط پایه با پارامترهای ثابت. بهینه‌ساز \lr{SLSQP} پنج پارامتر $(\mu, \nu, p_s, p_d, p_v)$ را به‌ازای هر فاصله تنظیم می‌کند.}
\label{fig:04_auto_opt}
\end{figure}

مقایسه با شکل~\ref{fig:01_baseline} نشان می‌دهد که بهینه‌سازی خودکار، نرخ کلید را به‌ویژه در فواصل میانی (۲۰ تا ۸۰ کیلومتر) بهبود می‌بخشد. در فواصل بسیار کوتاه، پارامترهای ثابت پیش‌فرض از قبل نزدیک به بهینه هستند و سود بهینه‌سازی ناچیز است. در فواصل بسیار طولانی، بهینه‌سازی می‌تواند برد امن را چند کیلومتر افزایش دهد اما نرخ نهایی همچنان به صفر میل می‌کند.

% ---------------------------------------------------------------------
\chapter{مقایسه‌های سیستماتیک}
\label{ch:results_comparisons}

\section{مقایسهٔ پروتکل‌ها}
شکل~\ref{fig:05_protocols} چهار پروتکل \lr{BB84-Decoy}، \lr{B92}، \lr{MDI-QKD} و \lr{Redundant Transmission} را در شرایط یکسان مقایسه می‌کند. پروتکل \lr{BB84-Decoy} در فواصل کوتاه تا متوسط بهترین عملکرد را دارد؛ \lr{MDI-QKD} به‌دلیل نیاز به آشکارسازی وسط مسیر، نرخ پایین‌تری دارد اما در برابر حملات آشکارساز مقاوم است و در فواصل طولانی نسبتاً پایدارتر عمل می‌کند. \lr{B92} ساده‌ترین پروتکل است اما نرخ پایین آن، کاربردش را در سناریوهای عملی محدود می‌سازد.

\begin{figure}[htbp]
\centering
\safeimg[width=0.85\textwidth]{figures/05_protocols_comparison.png}
\caption{مقایسهٔ نرخ کلید امن چهار پروتکل \lr{BB84-Decoy}، \lr{B92}، \lr{MDI-QKD} و \lr{Redundant Transmission} در شرایط فیبر یکسان.}
\label{fig:05_protocols}
\end{figure}

\textbf{یادداشت:} اگر نمودار فوق به‌صورت چهار فایل جداگانه ذخیره شده است، می‌توانید به‌جای شکل واحد بالا، چهار \lr{subfigure} کنار هم قرار دهید. در آن صورت، نام فایل‌ها به‌ترتیب \lr{\texttt{05\_protocol\_BB84DecoyProtocol.csv.png}}، \lr{\texttt{05\_protocol\_B92Protocol.csv.png}}، \lr{\texttt{05\_protocol\_MDIQKDProtocol.csv.png}} و \lr{\texttt{05\_protocol\_RedundantTransmissionProtocol.csv.png}} خواهند بود.

\section{مقایسهٔ اثبات‌های امنیتی}
شکل~\ref{fig:06_proofs} چهار اثبات امنیت \lr{Lim2014}، \lr{Ma2005}، \lr{Wang2005} و \lr{Tight} را روی همان پروتکل \lr{BB84-Decoy} مقایسه می‌کند. اثبات \lr{Lim2014} پیش‌فرض چارچوب است و تعادل خوبی بین سخت‌گیری و کارایی برقرار می‌کند. اثبات \lr{Tight} نرخ بالاتری تولید می‌کند اما فرضیات قوی‌تری دارد. اثبات‌های \lr{Ma2005} و \lr{Wang2005} برای سناریوهای خاص (مثل منابع نوری غیرایده‌آل) طراحی شده‌اند و در حالت استاندارد ممکن است نرخ پایین‌تری ارائه دهند.

\begin{figure}[htbp]
\centering
\safeimg[width=0.85\textwidth]{figures/06_proofs_comparison.png}
\caption{مقایسهٔ نرخ کلید امن بر اساس چهار اثبات امنیت مختلف برای پروتکل \lr{BB84-Decoy}.}
\label{fig:06_proofs}
\end{figure}

تفاوت بین این چهار منحنی، حاشیهٔ امنیت (\lr{security margin}) هر اثبات را کمّی می‌کند؛ هرچه اثبات سخت‌گیرانه‌تر باشد، حاشیه بزرگ‌تر و نرخ نهایی پایین‌تر است. انتخاب اثبات در عمل به مدل تهدید و میزان اطمینان به منبع نوری و آشکارساز بستگی دارد.

\section{مقایسهٔ آشکارسازها}
شکل~\ref{fig:07_detectors} سه نوع آشکارساز \lr{SPD}، \lr{SNSPD} و \lr{PNRD} را مقایسه می‌کند. آشکارساز \lr{SNSPD} به‌دلیل نرخ شمارش تاریک بسیار پایین ($\sim 10^{-8}$) و راندمان بالا ($\sim 0.5$)، بهترین عملکرد را در فواصل طولانی دارد. آشکارساز \lr{SPD} راندمان پایین‌تری دارد اما هزینه و پیچیدگی کمتری داشته و برای فواصل کوتاه مناسب است. آشکارساز \lr{PNRD} (آشکارساز شمارش فوتون متمایز) در سناریوهایی که نیاز به تشخیص دقیق تعداد فوتون‌ها است، کاربرد دارد.

\begin{figure}[htbp]
\centering
\safeimg[width=0.85\textwidth]{figures/07_detectors_comparison.png}
\caption{مقایسهٔ نرخ کلید امن با سه نوع آشکارساز \lr{SPD}، \lr{SNSPD} و \lr{PNRD}.}
\label{fig:07_detectors}
\end{figure}

این مقایسه نشان می‌دهد که ارتقای آشکارساز از \lr{SPD} به \lr{SNSPD} می‌تواند برد امن را تا دو برابر افزایش دهد، اما هزینه و نیاز به سرمایش کرایوژنیک ملاحظات عملی مهمی هستند که باید در طراحی سیستم لحاظ شوند.

\section{مقایسهٔ روش‌های بازهٔ اطمینان}
شکل~\ref{fig:08_confidence} سه روش \lr{Gaussian}، \lr{Clopper-Pearson} و \lr{Hoeffding} را برای محاسبهٔ بازهٔ اطمینان \lr{QBER} مقایسه می‌کند. روش \lr{Clopper-Pearson} سخت‌گیرانه‌ترین (محافظه‌کارانه‌ترین) روش است و بازه‌ای گسترده‌تر تولید می‌کند که منجر به کاهش نرخ کلید نهایی می‌شود. \lr{Gaussian} روش تقریبی است که برای نمونه‌های بزرگ بهینه است و نرخ بالاتری می‌دهد. \lr{Hoeffding} حد متوسطی ارائه می‌کند که بدون فرض توزیع خاص، همواره معتبر است.

\begin{figure}[htbp]
\centering
\safeimg[width=0.85\textwidth]{figures/08_confidence_comparison.png}
\caption{تأثیر روش محاسبهٔ بازهٔ اطمینان بر نرخ کلید امن. \lr{Clopper-Pearson} محافظه‌کارانه‌ترین و \lr{Gaussian} کارآمدترین روش است.}
\label{fig:08_confidence}
\end{figure}

انتخاب روش به تعداد پالس و سطح اطمینان مورد نظر بستگی دارد؛ برای $N < 10^5$، روش \lr{Clopper-Pearson} به‌دلیل دقت بهتر در نمونه‌های کوچک توصیه می‌شود و برای $N \geq 10^6$، تقریب \lr{Gaussian} کافی است و نرخ کارآمدتری به دست می‌دهد.

% ---------------------------------------------------------------------
\chapter{تحلیل حساسیت}
\label{ch:results_sensitivity}

\section{حساسیت به QBER ذاتی و \lr{Misalignment}}
شکل‌های~\ref{fig:10a_sensitivity} و~\ref{fig:10b_sensitivity} تأثیر هم‌زمان \lr{QBER} ذاتی آشکارساز و جابه‌جایی زاویه‌ای پایه‌ها (\lr{misalignment}) بر نرخ کلید امن را نشان می‌دهند. این دو پارامتر مهم‌ترین منابع خطا در سیستم‌های \lr{QKD} عملی هستند و افزایش هر کدام به‌طور مستقیم \lr{QBER} کلی را بالا برده و نرخ کلید امن را کاهش می‌دهد. در مقادیر بالا (مثلاً \lr{\texttt{qber\_intrinsic=0.02}} و \lr{\texttt{misalignment=0.01}})، حتی در فواصل کوتاه نیز نرخ کلید به‌طور قابل توجهی افت می‌کند.

\begin{figure}[htbp]
\centering
\safeimg[width=0.95\textwidth]{figures/10a_qber_misalignment_sweep.csv.png}
\caption{تحلیل حساسیت نرخ کلید امن به \lr{QBER} ذاتی آشکارساز.}
\label{fig:10a_sensitivity}
\end{figure}

\begin{figure}[htbp]
\centering
\safeimg[width=0.95\textwidth]{figures/10b_qber_misalignment_sweep.csv.png}
\caption{تحلیل حساسیت نرخ کلید امن به عدم انطباق پایه‌ها (\lr{misalignment}).}
\label{fig:10b_sensitivity}
\end{figure}

این تحلیل حساسیت دو پیامد عملی دارد: نخست، اهمیت کالیبراسیون دقیق پایه‌ها در زمان نصب سیستم \lr{QKD}؛ چرا که \lr{misalignment} حتی در حد چند درجه می‌تواند نرخ کلید را تا ۵۰ درصد کاهش دهد. دوم، انتخاب آشکارساز با \lr{QBER} ذاتی پایین (مانند \lr{SNSPD} در مقایسه با \lr{SPD}) تأثیر مستقیمی بر عملکرد نهایی سیستم دارد.

